% ============================================================
% Section 2: Background
% ============================================================
\section{Background}
\label{sec:background}

\subsection{XOR Arbiter PUF}

An Arbiter PUF (APUF) is a delay-based strong PUF that converts manufacturing
variation in circuit propagation delays into a binary response.  A challenge
$\mathbf{c} \in \{0, 1\}^{n}$ selects one of $2^n$ paths through a chain of
$n$ switching stages, each containing two multiplexers that pass or
cross-switch the signal depending on the challenge bit.  The arrival-time
difference at the two latches determines the response:
\[
r = \operatorname{sgn}\big(\Delta(\mathbf{c})\big) \in \{-1, +1\},
\]
where $\Delta(\mathbf{c})$ is the cumulative delay difference.  This is
conveniently modeled as a linear threshold function~\cite{ruhrmair2010modeling}.
The challenge $\mathbf{c}$ is first transformed into a feature vector
$\boldsymbol{\phi} \in \{-1, +1\}^{n+1}$ via the parity encoding:
\[
\phi_i = \prod_{j=i}^{n} (2c_j - 1), \quad i = 1, \ldots, n,
\quad \phi_{n+1} = 1.
\]
The response is then
\[
r = \operatorname{sgn}\big(\langle \mathbf{w}, \boldsymbol{\phi} \rangle\big),
\]
where $\mathbf{w} \in \mathbb{R}^{n+1}$ collects the per-stage delay
differences, distributed as $\mathbf{w} \sim \mathcal{N}(0, \alpha^2)$ with
$\alpha \ll 1$.

A $k$-XOR APUF combines $k$ independent APUFs by XORing their responses:
\[
r_{\text{XOR}} = \bigoplus_{j=1}^{k} r_j
= \prod_{j=1}^{k} \operatorname{sgn}\big(\langle \mathbf{w}_j,
\boldsymbol{\phi} \rangle\big).
\]
The XOR operation destroys the linear relationship between challenges and
responses.  The decision boundary is the product of $k$ hyperplanes, creating
a manifold whose complexity grows exponentially with $k$.  An adversary observing a pair $(\mathbf{c}, r_{\text{XOR}})$ learns only
one bit that depends on all $k$ weight vectors simultaneously; any single
APUF's contribution is masked by the others.

The conventional security claim is that a $k$-XOR APUF requires $O(2^k)$ CRPs
to model~\cite{ruhrmair2010modeling}.  This exponential scaling has made XOR
APUFs the de facto standard for security-critical applications, despite
growing evidence that practical attack thresholds fall below this bound.

\subsection{ML-Based Modeling Attacks}

Modeling attacks treat the PUF as an unknown function
$f: \mathcal{C} \to \{-1, +1\}$ and attempt to approximate it from observed
CRPs.  For single APUFs ($k=1$), the linear threshold structure makes
logistic regression with polynomial features highly effective, routinely
achieving $> 95\%$ accuracy with tens of thousands of
CRPs~\cite{ruhrmair2013}.

For $k \geq 2$, the XOR operation renders the response linearly inseparable,
requiring nonlinear models; multi-layer perceptrons (MLPs) are the standard
tool.  A typical MLP for XOR PUF modeling has 3--4 hidden layers with Tanh
activation and a Sigmoid output, trained with binary cross-entropy
loss~\cite{santikellur2023}.  Its hidden units approximate the multiplicative
XOR interaction, with their number typically chosen in proportion to $2^{k-1}$
to accommodate the $2^k$ possible XOR configurations.

Recent advances include:
\begin{itemize}
\item \textbf{Evolution strategies.} Becker~\cite{becker2015cia}
demonstrated CMA-ES attacks on 5-XOR and 6-XOR APUFs; the CalyPSO
framework~\cite{calypso2024} extended this to particle swarm optimization,
reporting success on 8-XOR with approximately 8M CRPs.

\item \textbf{Deep learning.} Fei~et~al.~\cite{fei2024mope} proposed the
MoPE and MMoPE frameworks, which adjust model complexity to PUF complexity
and report success on up to 7-XOR, showing that architectural matching
improves attack efficiency.

\item \textbf{Chosen-challenge attacks.} Sayadi~et~al.~\cite{sayadi2025}
showed that selectively choosing challenge patterns exploits systematic biases
in the PUF response distribution, reducing the CRP requirement and reaching
higher XOR configurations than random-challenge attacks.
\end{itemize}

Despite these advances, prior work has focused on \emph{whether} a given XOR
configuration can be attacked, reporting a single success point or best
accuracy; the \emph{probability} of attack success as a function of CRP
budget has not been systematically characterized, the gap this work
addresses.

\subsection{Split Attack and Interpose PUF}

The split attack, introduced by Wisiol~et~al.~\cite{wisiol2020splitting},
exploits the structural decomposition of XOR APUFs: rather than learning the
$k$-XOR function directly, it trains separate sub-models on subsets of APUFs
and combines their outputs probabilistically.  For a $k$-XOR
APUF, the response can be expressed as the XOR of two halves:
\[
r = r_{\text{up}} \oplus r_{\text{down}},\quad
r_{\text{up}} = \bigoplus_{j=1}^{\lfloor k/2 \rfloor} r_j,\quad
r_{\text{down}} = \bigoplus_{j=\lfloor k/2 \rfloor+1}^{k} r_j.
\]
The split attack models $P(r_{\text{up}}=1 \mid \mathbf{c})$ and
$P(r_{\text{down}}=1 \mid \mathbf{c})$ with two smaller MLPs, then combines
them as $P(r=1) = p_{\text{up}}(1-p_{\text{down}}) + p_{\text{down}}(1-p_{\text{up}})$.
This decomposition reduces the effective XOR complexity by a factor of
roughly 2, at the cost of training two models.  The NNModel\_fc variant uses a
single MLP with two output heads, achieving identical accuracy with
approximately $7.7\times$ fewer parameters than the standard approach
(\S\ref{sec:success_curves}).

The Interpose PUF (iPUF)~\cite{nguyen2019} was proposed as a
split-attack-resistant alternative.  An iPUF has a two-layer structure: an
upper layer of $k_{\text{up}}$ XORed APUFs produces a response $r_{\text{up}}$,
which is inserted into the challenge of a lower layer of $k_{\text{down}}$
XORed APUFs at a fixed bit position.
The security claim is that a $(k_{\text{up}}, k_{\text{down}})$-iPUF provides
security equivalent to a $(k_{\text{up}} + k_{\text{down}})$-XOR APUF, because
the attacker cannot separate the upper and lower layers.  Wisiol~et~al. refuted
this claim by adapting the split attack to iPUFs: the upper layer is attacked
independently and the lower layer conditioned on the predicted upper response.
Our results (\S\ref{sec:ipuf}) comprehensively confirm this vulnerability
across multiple configurations, showing that iPUF security is substantially
lower than claimed.
